% ======================================================================
%  08_outlook.tex —— 第 8 讲 收获与边界：证明了什么，没证明什么
%  依据：Kapustin--Witten 2006 \S12（印刷 pp.207--221）与全文边界
% ======================================================================

\section{收获与边界：证明了什么，没证明什么}
\label{sec:outlook}

\begin{goals}
\begin{itemize}
\item 看懂 \S12 的检查表：三类膜 $(B,B,B)$、$(B,A,A)$、$(A,B,A)$ 如何在四维规范理论里被边界条件造出来，为什么每一类都保两个拓扑超荷。
\item 能把这份推导的收获与缺口分开说：哪些是算出来的，哪些是「物理上应该」，哪些根本没做。
\item 知道 2024 年前后纯代数几何的证明与这条物理路线是什么关系，以及接下来该读什么。
\end{itemize}
\end{goals}

前十讲我们把一条链从四维拉到了二维：$\mathcal N=4$ 超对称 Yang--Mills 扭化、紧化到 $C$ 上、得到以 $\MH$ 为靶的 sigma 模型，再把 Wilson/'t~Hooft 线算子翻译成膜上的算子，把 't~Hooft 算子认成 Hecke 算子，最后用 c.c.\ 膜把 A-膜翻成 D-模（\cite{kw06}）。这一讲做三件事：把膜放回四维（\S12），清点这份推导的边界，再给出继续读下去的路线。

\subsection{回到四维：\S12 是一张检查表}

为什么最后还要回四维？因为几何 Langlands 的 Hecke 算子定义在 $\mathrm{Bun}_G(C)$ 上——所有 $G$-丛构成的栈，而不只是稳定丛的模空间 $\MC$；二维 sigma 模型的靶只有 $\MH$ 的光滑部分，在那里定义不了整体对象（\src{B p.184}）。原文把困难转成规范理论的类比：栈的规范理论对应物是所有联络的空间，于是「膜在二维定义不了的地方」，可以回到四维用边界条件定义。\S12 的目标原话就是：说明重要的膜确实可以在四维规范理论中定义（\src{B p.207}）。

四维的设定是 $M=\RR^2_+\times C$，边界在 $x^1=0$；一张膜就是边界上的一个超对称边界条件。做法是固定的四步：选一个保超对称的对合 $T$，猜费米子边界条件 $T\lambda|=\pm\lambda|$，用超对称变分 $\delta\lambda=\tfrac12\Gamma^{IJ}F_{IJ}\epsilon$ 推出玻色子边界条件，最后验证超流法向分量消失、并检查存在超对称态（\src{B pp.208--210}）。每一步都是线性条件，可以逐条核对，所以 \S12 读起来像清单而不像新理论——它不产出新结论，只验证前面用到的膜确有四维来源。

三类膜是清单的主角。$(B,B,B)$ 型是 $\MH(\Ldual G,C)$ 上的 zerobrane：一个点，在三个复结构里都是 B-膜。它的镜像是 Hitchin 纤维——对 $I$ 复、对 $J$ 与 $K$ Lagrange，故是 $(B,A,A)$ 型。c.c.\ 膜是 $(A,B,A)$ 型，它的镜像也是（\src{B p.207}）。这里的 $I,J,K$ 就是第 1 讲（原文 §4.1）介绍的三个复结构：一张膜在某个复结构中算 A 还是 B，看它保住的是那个结构的 A-型还是 B-型拓扑超荷；同一张膜相对三个结构可以有不同的类型，这正是三元组记号的读法。

为什么每一类都恰好保两个拓扑超荷？扭化后的理论里有两个标量超荷 $Q_l,Q_r$（它们的任意非零线性组合 $Q=uQ_l+vQ_r$ 都可以充当拓扑超荷），而边界条件一旦保住 $Q$，也自动保住它的共轭 $Q^\dagger$；两个拓扑超荷加上各自的共轭，共四个超对称——是 $\mathcal N=4$ 的 16 个中的四分之一，所以这类膜叫 $1/4$ BPS 膜；从只有一半超对称的二维 sigma 模型看，它们反而是 $1/2$ BPS（\src{B p.208}）。三种类型之所以都保同样两个超荷，是因为它们在不同复结构中的 A/B 类型，最终都化为 $Q_l,Q_r$ 的线性组合。

三类膜的四维实现可以各记一个标志性条件。$(B,B,B)$：取 $T=\Gamma^{04}$，玻色条件是 $A_0+A_4=0$、其余 $A_I$ 取 Dirichlet，而超对称态条件化成 Hitchin 方程，于是边界值给出 $\MH$ 上的一个点（\src{B pp.211--213}）。$(B,A,A)$：对 $(B,B,B)$ 的条件作 S-对偶，$T$ 变成 $\Gamma^1\Gamma^2\Gamma^3\Gamma^4$，剩下的条件等价于复 Hitchin 方程，膜就卷在 Hitchin 纤维上（\src{B p.216}）。$(A,B,A)$：混合边界条件一部分 Dirichlet、一部分 Neumann，把降维后的条件与 sigma 模型的混合边界条件对比，可读出目标空间规范场恰是 $f=\omega_J$——正是 \S11 里 c.c.\ 膜的曲率（\src{B p.220}）。这一步是全讲的闭环：\S11 用 c.c.\ 膜造出 D-模，\S12 又从四维边界条件把同一张膜造回来。

\begin{table}[htbp]
\centering
\footnotesize
\begin{tabular}{@{}l p{0.19\textwidth} c p{0.23\textwidth} p{0.20\textwidth}@{}}
\toprule
膜类型 & 在 $\MH$ 上的几何 & $(I,J,K)$ & 保哪两个超荷 & 四维实现要点 \\
\midrule
$(B,B,B)$ & zerobrane：支撑在一个点 & $B,B,B$ &
两个拓扑超荷 $Q_l,Q_r$；三个结构中的 B-型 BRST 都是它们的线性组合 &
$T=\Gamma^{04}$；$A_0+A_4=0$，其余 Dirichlet \\
$(B,A,A)$ & Hitchin 纤维：固定 $\varphi_z$ 的特征多项式 & $B,A,A$ &
同样是 $Q_l,Q_r$；$I$ 中取 B-型、$J,K$ 中取 A-型的组合 &
$T'=\Gamma^{1}\Gamma^{2}\Gamma^{3}\Gamma^{4}$；复 Hitchin 方程 \\
$(A,B,A)$ & space-filling：全 $\MH$，带曲率 $F\propto\omega_J$ & $A,B,A$ &
同样是 $Q_l,Q_r$；$I,K$ 中取 A-型、$J$ 中取 B-型的组合 &
$\hat T$ 的混合边界条件；$f=\omega_J$ 认出 c.c.\ 膜 \\
\bottomrule
\end{tabular}
\caption{三类关键膜的四维对照。三元组记录的是同一张膜相对三个复结构 $I,J,K$ 的 A/B 类型，不是三种不同的膜理论；三者都保两个拓扑超荷（连共轭共四个超对称），所以都是 $1/4$ BPS（\src{B p.207}；\src{B p.208}）。}
\label{tab:three-branes}
\end{table}

清单也有它自己的诚实之处：提升不唯一。原文明确说，某些情形会给出多个不同的四维边界条件，它们在降到二维后、至少在奇点之外看起来等价；$(B,B,B)$ 的第二种构造在平坦时空中保八个超荷而不是四个，两种构造只在奇点附近才可能显出差别（\src{B pp.213--214}）。因此 \S12 检查的是「存在性」，不是「唯一性」。

\subsection{这份推导给了什么：一个图像、一套对应、大量证据}

把整条链压成一句话：\kw{} 用四维 $\mathcal N=4$ 理论的 S-对偶实现了几何 Langlands 的两边——左边是局部系统（magnetic eigenbrane），右边是 D-模（经 c.c.\ 膜翻译），中间的桥是 S-对偶在二维变成镜像对称（\cite{kw06}）。它给出的收获分三层。第一层是\textbf{物理图像}：几何 Langlands 不是孤立猜想，而是电磁对偶在特定紧化下的表现。第二层是\textbf{对应关系}：对象层（eigenbrane $\leftrightarrow$ Hecke 特征层）与算子层（'t~Hooft $\leftrightarrow$ Hecke）各有一套明确的字典。第三层是\textbf{证据}：$U(1)$ 情形可以完整算完，三类膜的类型表自洽，\S12 用四维边界条件独立造出 \S11 用到的膜，边界算子的 Q-上同调又与开弦代数对上。

但它是推导，不是证明。原文用 will show、roughly speaking、we claim、we argue 这些词，\S11 的翻译是单向构造，「范畴等价」只是转引（\src{B p.184}）。三处最明显的软肋是：其一，翻译与 Hecke 作用的\textbf{交换性}没有逐条证明——若「先翻译再作用」与「先作用再翻译」不一致，得到的 D-模未必是 Hecke 特征层；其二，层的完全粘接只在 $T^*V$ 型开集上成立，精确理论里有收敛问题（\src{B p.190}）；其三，非整数 $\Psi$ 处的 $f(\Psi)$ 只由微扰论证确定（\src{B p.207}）。

所以我们建议这样记：\kw{} 给的是\textbf{为什么}，不是\textbf{是什么的证明}。物理推导的可信度来自独立路径给同一答案，而不是来自某一步的形式化；读它时要能指出每一步的证据强度，这正是 \S12 存在的意义——它是这份推导的诚实度指标。

\subsection{边界在哪里：非分歧、二维 CFT、以及 2024 年的证明}

原文自己标出了两条边界。第一条：只做非分歧（unramified）情形，即只考虑 $C$ 上没有分歧点的局部系统，分歧与野分歧完全没碰（\src{B p.5}）；分歧情形要在 $\mathrm{Bun}_{G,y}$ 上加 level structure、引入仿射 Hecke 代数，是 Gukov--Witten 的主题（\cite{gw06}）。第二条：没有说清二维共形场论在几何 Langlands 里到底起什么作用（\src{B p.5}）。这不是小遗漏：Frenkel 那一派恰恰以二维 CFT 为主工具——顶点代数、共形块、临界 level $k=-h^\vee$ 的仿射代数——从另一头构造 Hecke 特征层（\cite{fr09}）。两条路线对象相同、入口相反：\kw{} 解释这条对应为什么存在，Frenkel 的路线构造它是什么。

还有两个技术边界值得记住：D-模本身必须做在 $\mathrm{Bun}_G(C)$ 这个栈上，只在稳定丛上做不闭合（\src{B p.185}）；某些陈述对 $C$ 的 spin 结构敏感，$\Psi=h$ 的普通 D-模陈述也要据此修正（\src{B p.207}）。

那么数学上到底证明了没有？2006 年时没有。2024 年前后，Gaitsgory 与合作者用五篇系列预印本给出了特征 $0$ 代数闭域上光滑射影曲线的几何 Langlands 的证明（D-模、非分歧情形）。要把这件事的位置说准：它是\textbf{数学定理}，走的是纯代数几何与范畴化的路线，既不用 S-对偶，也不用膜；它不验证 \kw{} 的物理论证，\kw{} 也没有预言它的方法。两者的关系是分工：物理路线提供直觉、猜想与「为什么应该对」，数学路线提供定理与「确实对」。有意思的是，定理的适用范围（非分歧、光滑射影曲线）与 \S12 检查表的适用范围恰好一致——这不是巧合，而是 \kw{} 一开始就把自己限制在能做干净的那一块。

\subsection{接下来读什么：文献与三条建议}

如果只再读一份，读 Witten 2008 的十五页短文（\cite{w08}）：它把「紧化 $+$ 镜像对称」这一段重写成 Hitchin 纤维化上的镜像对称，是整条链最短的物理重述。想补二维 CFT 那一侧，读 Frenkel 的短综述（\cite{fr09}），它会告诉你 Hecke 特征层在 CFT 里是怎么被造出来的。想要更全的对照，可用 Ben-Zvi 的综述（\cite{bz04}）与 Hitchin 系统的材料（\cite{ht03}）。要往分歧情形走，读 Gukov--Witten（\cite{gw06}）；想把框架推到三维与面算子范畴，读 Setter 的博士论文（\cite{set13}）——注意它本身既不证明几何 Langlands，也不定义 Hecke 作用，读它是学技术而不是找结论。

给读者三条建议。第一，\textbf{先读第 5、6 讲再回原文 \S11}：c.c.\ 膜如果没有 eigenbrane 与 Hecke 算子的动机，看起来只是一张奇怪的膜；有了动机，才知道它为什么必须空间充满。第二，\textbf{\S12 可以当检查表跳读}：只需记住三类膜、各自的一句话来源，以及「提升不唯一」这一条；四步法的细节等要用时再回查。第三，\textbf{卡住时用 $U(1)$ 例子校准}：\S11.2 末的 $U(1)$ 计算是全文唯一能从头算到尾的例子，把它走一遍，$\lambda$-连接、平坦联络、Poincaré 对偶这些词就都有了落点（\src{B pp.199--201}）；再卡，就查解经本对应页段的骨血层。

\begin{keybox}
\kw{} 交给几何 Langlands 的是一份物理推导：一个 S-对偶图像、一套膜与算子的字典、以及 \S11 与 \S12 两方向互相对上的大量证据；它没给证明，只做非分歧情形，也不解释二维 CFT 的作用——真正的定理要等 2024 年前后纯代数几何的路线。
\end{keybox}

\begin{follow}
\begin{itemize}
\item 先读 \S12 的开头三页（\src{B pp.207--210}）：目标、三类膜、为什么保两个拓扑超荷、四步法，全在这里。
\item 再挑一类细读：$(B,B,B)$ 在 \src{B pp.211--215}，$(A,B,A)$ 与 c.c.\ 膜的识别在 \src{B pp.217--221}；后者与 \S11 的接口最值得看。
\item 边界部分读原文 \S1 的结尾（\src{B p.5}）：作者自陈的两个 gap（非分歧、二维 CFT），是理解这份推导定位的关键。
\item 后续文献按顺序：\cite{w08} 十五页 $>$ \cite{fr09} 三十二页 $>$ \cite{bz04} 综述 $>$ \cite{gw06} 分歧 $>$ \cite{set13} 博士论文。
\end{itemize}
\end{follow}

\begin{checkbox}
\begin{itemize}
\item 三类膜各对应 $(I,J,K)$ 中的哪种 A/B 类型？为什么同一张膜可以在三个复结构里有不同身份？
\item 「保两个拓扑超荷」为什么等于「共四个超对称」？$1/4$ BPS 与 $1/2$ BPS 两个说法分别从哪个理论的角度看？
\item \S12 为什么叫检查表而不是新理论？举一个它检查出「提升不唯一」的例子。
\item 原文自陈的两个 gap 是什么？Frenkel 的路线如何从另一头绕开第二个 gap？
\item 2024 年的证明与 \kw{} 的物理路线是什么关系？为什么不能说「物理证明了几何 Langlands」？
\end{itemize}
\end{checkbox}
